%%=============================================================================
%% CAPITULO 5 -- CONSIDERACOES FINAIS
%% Esqueleto provisorio (dissertacao dummy): texto a ser desenvolvido.
%%=============================================================================

\section{Considerações finais}

A pesquisa partiu de um ativo digital público concreto, as 760 videoaulas do
Preparatório ENEM do Canal Educação da SEDUC-PI, e perguntou em que medida
alternativas de classificação automatizada conseguem atribuir-lhes as
habilidades da matriz de referência do ENEM. Os resultados indicam a
viabilidade técnica da recomendação baseada em conteúdo em contexto de início
frio: o artefato com embeddings BGE-M3 supera o método lexical sem inteligência
artificial na concordância com a atribuição oficial do INEP, mas fica aquém do
modelo generativo.

As limitações devem ser declaradas sem atenuação. Item de prova e videoaula não
são o mesmo objeto, de modo que a concordância medida sobre itens do ENEM não é
medida de acerto do artefato em sua função; a aferição sobre as próprias
videoaulas exige padrão-ouro construído por professores especialistas, que
integra a proposta de continuidade da pesquisa, junto com a aferição de
utilidade percebida pelos docentes. A causa da fragilidade específica em
Linguagens e Códigos permanece questão aberta.

Do ponto de vista da gestão pública, a classificação do acervo converte-se em
insumo de governança do próprio ativo, ao expor as lacunas de cobertura
curricular que orientam a priorização da produção de novos conteúdos. A
recomendação para a rede que priorize autonomia e custo é a representação
vetorial local, com o alvo construído a partir dos itens publicados pelo exame,
por ser a melhor concordância obtida sem custo por requisição e sem transmissão
de conteúdo a terceiro.
